\clearpage
\section*{Appendix R\quad Computational Receipts}\label{app:receipts}
\addcontentsline{toc}{section}{Appendix R: Computational Receipts}
\small\noindent Each computation marked \rcptmarker\ in the text is verified by a runnable receipt in the \texttt{receipts/} directory that ships with this work; status \textsf{OK} = verified (traced, run, output matches the claim). This appendix is generated from the receipt index, not maintained by hand.\par\medskip
\begingroup\renewcommand{\arraystretch}{1.25}
\begin{description}
\item[\label{rcpt:P10_minisuperspace_friedmann}\texttt{P10\_minisuperspace\_friedmann}]\hfill\textsf{[OK]}\\
\textit{sec:deparam} \ (deparametrized minisuperspace Hphys=(2pi/3)p\_a\textasciicircum{}2/a-(L/8pi)a\textasciicircum{}3 generates Friedmann (adot/a)\textasciicircum{}2=(8pi/3)rho\_d/a\textasciicircum{}3+L/3 via Hamilton's equations). 
\emph{Computes:} adot=4pi p\_a/3a; Hphys autonomous=conserved; eliminate p\_a -> Friedmann; dust a\textasciicircum{}-3; Lambda-sign control
 \emph{Bound:} the true-Hamiltonian structure is non-empty
\ \emph{Run:} \texttt{python3 receipts/P10\_canonical\_time/P10\_minisuperspace\_friedmann.py}
\item[\label{rcpt:P10_graviton_lift}\texttt{P10\_graviton\_lift}]\hfill\textsf{[OK]}\\
\textit{sec:lock} \ (closed-S\textasciicircum{}3 graviton lift: tensor-harmonic action S\_n=1/2 int dT a\textasciicircum{}3[phidot\textasciicircum{}2-(mu\textasciicircum{}2/a\textasciicircum{}2)phi\textasciicircum{}2] Legendre-transforms to Hphys\_n=pi\textasciicircum{}2/2a\textasciicircum{}3+1/2 a mu\textasciicircum{}2 phi\textasciicircum{}2; mode eq phiddot+3(adot/a)phidot+(mu\textasciicircum{}2/a\textasciicircum{}2)phi=0). 
\emph{Computes:} pi=a\textasciicircum{}3 phidot; Legendre H; EL mode eq; Hamilton eqns combine to it; flat-a oscillator control
 \emph{Bound:} unitary cosmic-time evolution of the graviton tower
\ \emph{Run:} \texttt{python3 receipts/P10\_canonical\_time/P10\_graviton\_lift.py}
\item[\label{rcpt:LIFT_adiabatic_correction}\texttt{LIFT\_adiabatic\_correction}]\hfill\textsf{[OK]}\\
\textit{eq:adiabatic-exponent} \ (THE ADIABATIC CORRECTION. Exact Euclidean suppression for mode n is exp(-int omega\_n ds), with omega\_n = mu\_n / a (P10 sec:lock) and a = \textbar{}r\textbar{} along the lift. omega DIVERGES at the branch point (a->0). Does the exponent converge? And is the evolution adiabatic (\textbar{}dw/ds\textbar{}/w\textasciicircum{}2 << 1)?). 
\emph{Computes:} runs clean; registered from its own docstring and a live run
 \emph{Bound:} description is the receipt's own wording, condensed; not independently re-derived here
\ \emph{Run:} \texttt{python3 receipts/P10\_canonical\_time/LIFT\_adiabatic\_correction.py}
\item[\label{rcpt:LIFT_gravitational_action}\texttt{LIFT\_gravitational\_action}]\hfill\textsf{[OK]}\\
\textit{eq:grav-action} \ (THE CALIBRATION. On-shell S = int p\_a da with p\_a = (3/4pi) a adot (from P10's H). Continuing tau = i s : S = -i S\_E with S\_E\textasciicircum{}grav = (3/4pi) int r r'\textasciicircum{}2 ds along the SAME trajectory. Compare with the contour action S\_E\textasciicircum{}c = int [ (1/2) r'\textasciicircum{}2 + f/2 ] ds = 1.489 (alpha=1).). 
\emph{Computes:} runs clean; registered from its own docstring and a live run
 \emph{Bound:} description is the receipt's own wording, condensed; not independently re-derived here
\ \emph{Run:} \texttt{python3 receipts/P10\_canonical\_time/LIFT\_gravitational\_action.py}
\item[\label{rcpt:LIFT_instanton_action}\texttt{LIFT\_instanton\_action}]\hfill\textsf{[OK]}\\
\textit{eq:euclidean-action} \ (FAMILY 10, final item: does the lift have a variational principle of its own? Claim: S\_E = int ds [ (1/2) r'\textasciicircum{}2 + V(r) ], V = f/2, with Euclidean first integral (1/2)r'\textasciicircum{}2 - V = -1/2. Test the closed form r(s) = -(2 M a\textasciicircum{}2)\textasciicircum{}\{1/3\} \textbar{}sin(3s/2a)\textbar{}\textasciicircum{}\{2/3\} against BOTH.). 
\emph{Computes:} runs clean; registered from its own docstring and a live run
 \emph{Bound:} description is the receipt's own wording, condensed; not independently re-derived here
\ \emph{Run:} \texttt{python3 receipts/P10\_canonical\_time/LIFT\_instanton\_action.py}
\item[\label{rcpt:P10_the_straddle_does_not_need_a_floor}\texttt{P10\_the\_straddle\_does\_not\_need\_a\_floor}]\hfill\textsf{[OK]}\\
\textit{sec:lock} \ (**THE STRADDLE DOES NOT NEED A FLOOR, AND THE PARAGRAPH'S TWO STATEMENTS ABOUT Gamma-hat ARE ABOUT DIFFERENT OPERATORS.** The coupled sector promotes the boundary coefficient from the c-number gamma <= 1/4 to an operator Gamma-hat on the tower. **The paragraph asserted it 'bounded below by gamma' where it builds the straddle, and listed 'whether it is bounded below' as OPEN four hundred characters later.** Both are right of different operators: the LEADING-ORDER Gamma-hat = gamma + c sum\_n pi\_n\textasciicircum{}2 is a sum of squares over a positive coefficient, hence spectrum in [gamma, oo); the COMPLETE one carries the cubic and higher self-interactions, which enter at the same inverse-square order and are not sign-definite. ** AND THE REPAIR COSTS THE ARGUMENT NOTHING, WHICH IS THE POINT: the direct-integral decomposition of -d\textasciicircum{}2/dx\textasciicircum{}2 + Gamma-hat/x\textasciicircum{}2 uses only that the spectrum MEETS BOTH SIDES of 3/4 -- limit point at or above, limit circle below -- and a spectral projection onto \{Gamma-hat < 3/4\} exists whether or not the spectrum has a floor.** Exhibited on a two-channel model in which the cubic takes one channel: the union is unbounded below and both spectral subspaces stay non-trivial.). 
\emph{Computes:} the leading-order infimum asserted equal to gamma and the operator asserted monotone in k; the complete model asserted to break the floor and asserted driven down at large momentum; both spectral subspaces asserted non-trivial in both cases
 \emph{Bound:} **bound: the model is a diagonal two-channel stand-in for the singular part of the interacting graviton Hamiltonian, not that Hamiltonian. It establishes what the decomposition NEEDS, not what the spectrum IS -- and the spectrum is exactly the frontier the paragraph names (`L-165`)**
\ \emph{Run:} \texttt{python3 receipts/P10\_canonical\_time/P10\_the\_straddle\_does\_not\_need\_a\_floor.py}
\item[\label{rcpt:P10_the_straddle_is_computed}\texttt{P10\_the\_straddle\_is\_computed}]\hfill\textsf{[ALL PASS]}\\
\textit{sec:lock} \ (THE STRADDLE AS A COMPUTED FACT: spec(Gamma-hat)=[gamma,inf) from the displayed form, and gamma<=1/4<3/4 puts spectrum strictly below the threshold while unboundedness puts it strictly above, so both sides are occupied for every admissible gamma; the 3/4 exhibited by solving the indicial equation rather than quoted). 
\emph{Computes:} r3803
 \emph{Bound:} 59
\ \emph{Run:} \texttt{python3 receipts/P10\_canonical\_time/P10\_the\_straddle\_is\_computed.py}
\item[\label{rcpt:P10_the_adiabatic_residual_at_low_n_is_bounded_by_the_towers_own_floor}\texttt{P10\_the\_adiabatic\_residual\_at\_low\_n\_is\_bounded\_by\_the\_towers\_own\_floor}]\hfill\textsf{[ALL PASS]}\\
\textit{sec:dissolution} \ (**THE ADIABATIC RESIDUAL AT LOW n IS BOUNDED BY THE TOWER'S OWN FLOOR, AND IT REACHES NOTHING OBSERVABLE.** The suppression falls monotonically with harmonic index (2.8e-4 at n=2, 5.9e-6 at n=3, \textasciitilde{}2 orders per step after), so the mixing the adiabatic projection neglects can only carry amplitude into MORE-suppressed modes; and the S\textasciicircum{}3 tensor tower has no mode below n=2, so the least-suppressed mode has no source beneath it. **The bound is therefore one-sided by construction rather than by estimate.** Worst-case transfer of the whole O(eps) fraction out of n=2 leaves that mode MORE suppressed and does not raise the ceiling. Ceiling on transmitted power = unmixed n=2 = **7.9e-8**, and below 1e-3 even on the naive exponent.). 
\emph{Computes:} r4231
 \emph{Bound:} 61
\ \emph{Run:} \texttt{python3 receipts/P10\_canonical\_time/P10\_the\_adiabatic\_residual\_at\_low\_n\_is\_bounded\_by\_the\_towers\_own\_floor.py}
\item[\label{rcpt:P10_gamma_hat_is_bounded_below}\texttt{P10\_gamma\_hat\_is\_bounded\_below}]\hfill\textsf{[OK]}\\
\textit{sec:lock} \ (**THE COMPLETE BOUNDARY OPERATOR IS BOUNDED BELOW, AND THE REASON IS THE DEPARAMETRIZATION ITSELF: THE SCALE FACTOR IS THE SUPERMETRIC'S ONE NEGATIVE DIRECTION.** P10 lists 'the spectrum of Gamma-hat, including whether it is bounded below' among what is open in the interacting tower; that clause is separable and is settled here. **The DeWitt form tr((g pi)\textasciicircum{}2) - lambda (tr g pi)\textasciicircum{}2 has exactly ONE negative eigenvalue on the round metric, realised by the identity direction --- the conformal one, which is the whole indefiniteness of superspace and the reason Wheeler-DeWitt is hyperbolic.** On g-TRACELESS momenta the trace term drops and the form is exactly tr((g pi)\textasciicircum{}2), which for L the Cholesky factor of g equals tr(S\textasciicircum{}2) with S = L\textasciicircum{}T pi L SYMMETRIC --- a sum of squares of real eigenvalues, zero only at pi = 0. **Positive definite for every non-degenerate spatial metric, verified as an IDENTITY on 400 random pairs rather than as a sample.** **So the cubic term's apparent unboundedness is a TRUNCATION ARTEFACT: pi\textasciicircum{}2(1 - lambda phi + ...) is the expansion of pi\textasciicircum{}2/(1 + lambda phi), and the full coefficient is positive wherever the spatial metric is non-degenerate.** The deparametrization that produces a true Hamiltonian is the same move that produces a bounded-below one --- both are the removal of the conformal direction, and P10 argues the first without noticing it had the second.). 
\emph{Computes:} the signature asserted one negative; the Cholesky identity and positivity asserted on 400 valid (g, pi) pairs; min Q asserted positive on three metrics including a 625:1 anisotropy; the truncation asserted to go negative where the full expression does not
 \emph{Bound:} **bound: positivity holds on NON-DEGENERATE spatial metrics. At the degenerate boundary of superspace the inverse metric blows up and the argument says nothing --- which is exactly the a -> 0 boundary, handled instead by the thermal condition. Two different jobs, neither substituting. NOT TOUCHED: the UV definition of the tower sums, and the closed-form nonlinear solution**
\ \emph{Run:} \texttt{python3 receipts/P10\_canonical\_time/P10\_gamma\_hat\_is\_bounded\_below.py}
\item[\label{rcpt:I9_the_adiabatic_parameter_diverges_where_the_action_integral_converges}\texttt{I9\_the\_adiabatic\_parameter\_diverges\_where\_the\_action\_integral\_converges}]\hfill\textsf{[OK]}\\
\textit{sec:lock} \ (I9 --- THE ADIABATIC PARAMETER DIVERGES WHERE THE ACTION INTEGRAL CONVERGES. Integrable-systems field bake, probe I9. On P10's own near-branch-point scaling abs(r) \textasciitilde{} s\textasciicircum{}(2/3), the adiabaticity parameter abs(dw/ds)/w\textasciicircum{}2 goes as s\textasciicircum{}(-1/3) and DIVERGES, so the adiabatic invariant E/omega is not conserved through the branch point; the suppression exponent int omega ds nevertheless converges as S\textasciicircum{}(1/3), because omega \textasciitilde{} s\textasciicircum{}(-2/3) is integrable. Two different powers, two different questions: the finiteness P10 verifies comes from the action integral and not from slow variation, so the correction is semiclassical -- which is what P15 calls the same object. Two adversarial controls, one for each half.). 
\emph{Computes:} runs clean; eight verdicts including controls that fire in opposite directions
 \emph{Bound:} COMPUTES: the scaling abs(r) \textasciitilde{} s\textasciicircum{}(2/3) is P10's own; A and mu carried symbolically, no value pinned
\ \emph{Run:} \texttt{python3 receipts/P10\_canonical\_time/I9\_the\_adiabatic\_parameter\_diverges\_where\_the\_action\_integral\_converges.py}
\item[\label{rcpt:P10_gamma_is_one_quarter_and_is_the_maximum}\texttt{P10\_gamma\_is\_one\_quarter\_and\_is\_the\_maximum}]\hfill\textsf{[OK]}\\
\textit{sec:lock} \ (gamma = 1/4 is the MAXIMUM of the natural ordering family, so (1,1) holds for every ordering). 
\emph{Computes:} reduces T\_s exactly, shows the first-derivative term vanishes, and gets Gamma(s) with its maximum; 11 checks
 \emph{Bound:} NOT-A-PAPER-CLAIM --- verification of a standing claim; node 17 withdrew its own caveat
\ \emph{Run:} \texttt{python3 receipts/P10\_canonical\_time/P10\_gamma\_is\_one\_quarter\_and\_is\_the\_maximum.py}
\item[\label{rcpt:D1_the_boundary_is_per_fibre_and_the_UV_is_over_fibres}\texttt{D1\_the\_boundary\_is\_per\_fibre\_and\_the\_UV\_is\_over\_fibres}]\hfill\textsf{[OK]}\\
\textit{sec:lock} \ (PO-6 re-weighted not closed: the boundary closure is PER-FIBRE and the UV is OVER-FIBRES, and the clause c54.129 answered is the one the argument does not need). 
\emph{Computes:} asserts the direct-integral structure at source, computes the per-fibre deficiency across the threshold, and checks the straddle under an unbounded-below operator (18 checks)
 \emph{Bound:} NOT-A-PAPER-CLAIM --- a bounded re-statement of a PROTECTED row; PO-6 stays open
\ \emph{Run:} \texttt{python3 receipts/L165\_interacting\_tower/D1\_the\_boundary\_is\_per\_fibre\_and\_the\_UV\_is\_over\_fibres.py}
\item[\label{rcpt:D2_the_UV_degree_is_quartic_and_the_IR_is_free}\texttt{D2\_the\_UV\_degree\_is\_quartic\_and\_the\_IR\_is\_free}]\hfill\textsf{[OK]}\\
\textit{sec:lock} \ (PO-6s UV clause measured: the tower sum diverges QUARTICALLY, the ordinary zero-point degree, and compactness buys the IR rather than the UV). 
\emph{Computes:} takes the tower and mode action from P10 at source, derives the linear growth of mu\_n and <pi\_n\textasciicircum{}2>, and tracks the partial sums against N\textasciicircum{}4/4 (13 checks)
 \emph{Bound:} NOT-A-PAPER-CLAIM --- a measurement of one PO-6 clause; PO-6 stays open
\ \emph{Run:} \texttt{python3 receipts/L165\_interacting\_tower/D2\_the\_UV\_degree\_is\_quartic\_and\_the\_IR\_is\_free.py}
\item[\label{rcpt:S1_the_subleading_log_divergence_is_gauss_bonnet_and_the_ledger_survives}\texttt{S1\_the\_subleading\_log\_divergence\_is\_gauss\_bonnet\_and\_the\_ledger\_survives}]\hfill\textsf{[OK]}\\
\textit{---} \ (PO-6 fixed-background half: the SUB-LEADING (log) one-loop divergence on the tower's own de Sitter slicing is the Gauss-Bonnet term (topological in D=4), so it forces NO curvature-squared coupling outside the one-constant ledger --- extending L-809's quartic result to the log order; on de Sitter Weyl\textasciicircum{}2=0, and the on-shell graviton divergence reduces to GB (t Hooft-Veltman/Christensen-Duff)). 
\emph{Computes:} verifies R\textasciicircum{}2/Ric\textasciicircum{}2/Riem\textasciicircum{}2=144/36/24 per alpha\textasciicircum{}4 (S50) and Weyl\textasciicircum{}2=0, that GB=24/alpha\textasciicircum{}4 is the topological Euler density, the on-shell R=4Lambda, and the ledger-survival conclusion (4 checks)
 \emph{Bound:} NOT-A-PAPER-CLAIM --- the fixed-background half of PO-6; extends L-809; does NOT touch the coupled-sector dark half (54's/definitional)
\ \emph{Run:} \texttt{python3 receipts/L816\_po6\_subleading\_survives/S1\_the\_subleading\_log\_divergence\_is\_gauss\_bonnet\_and\_the\_ledger\_survives.py}
\item[\label{rcpt:S1_the_one_constant_ledger_survives_on_the_running_layer_not_only_fixed_de_sitter}\texttt{S1\_the\_one\_constant\_ledger\_survives\_on\_the\_running\_layer\_not\_only\_fixed\_de\_sitter}]\hfill\textsf{[OK]}\\
\textit{sec:lock} \ (PO-6 RUNNING-LAYER half (the residue S50 named: does the one-constant basis survive on a background whose curvature RUNS?): on a(T)=sinh\textasciicircum{}\{2/3\}(3T/2alpha), where R runs (R->12/alpha\textasciicircum{}2 only asymptotically), the pure-graviton log divergence reduces to \{Gauss-Bonnet (topological) + Lambda-renorm + Newton/ell\_P-renorm + field-equation-proportional (field-redef removable)\} with NO irreducible curvature-squared coupling outside \{Lambda,G\} --- for ANY one-loop coefficients. Extends L-816 off fixed de Sitter, via background-independent algebraic identities (not R=const)). 
\emph{Computes:} confirms Weyl\textasciicircum{}2=0 and R runs on the sinh\textasciicircum{}\{2/3\} layer, the topological (53/90)E\_GB split, the Lambda=0 identity (7/20)Ric\textasciicircum{}2+(1/120)R\textasciicircum{}2=E\_mn f\textasciicircum{}mn exactly, the coefficient-independent reduction of a1 Riem\textasciicircum{}2+a2 Ric\textasciicircum{}2+a3 R\textasciicircum{}2 to \{E\_GB,Lam\textasciicircum{}2,LamR,EOM\} with remainder 0, and the vacuum-dS consistency with L-816 (6 checks)
 \emph{Bound:} NOT-A-PAPER-CLAIM --- the running-layer half of PO-6; F5 reserves the verdict; pure graviton (matter as bend), does NOT enter the quantised-background sector or claim one-loop finiteness
\ \emph{Run:} \texttt{python3 receipts/L818\_po6\_running\_layer/S1\_the\_one\_constant\_ledger\_survives\_on\_the\_running\_layer\_not\_only\_fixed\_de\_sitter.py}
\item[\label{rcpt:S1_the_ledgers_one_irreducible_survivor_is_the_conformal_coupling_that_turns_on_at_the_shear}\texttt{S1\_the\_ledgers\_one\_irreducible\_survivor\_is\_the\_conformal\_coupling\_that\_turns\_on\_at\_the\_shear}]\hfill\textsf{[OK]}\\
\textit{sec:lock} \ (PO-6: L-818's boundary LOCATED with a coefficient (answers 56's r2764 routing). The graviton log counterterm = (7/40)Weyl\textasciicircum{}2 + (149/360)GB + (1/8)R\textasciicircum{}2; GB is topological and R\textasciicircum{}2 reduces (L-818), so the ONE irreducible piece is the conformal (7/40)Weyl\textasciicircum{}2 --- it is 0 on the shear-free FRW layer (which is WHY L-818's remainder was exactly 0 and the ledger survived) and turns on at 2nd order in the shear, where Weyl\textasciicircum{}2=4sigma\textasciicircum{}2 (56 S10/r2743). The shear IS the TT-graviton tower = P10's open interacting tower; so L-818's caveat, r2743's shear and P10's open item are ONE boundary, now carrying (7/40)Weyl\textasciicircum{}2). 
\emph{Computes:} decomposes the graviton b\_4 into \{Weyl\textasciicircum{}2,GB,R\textasciicircum{}2\} (Weyl\textasciicircum{}2 coeff 7/40), shows only Weyl\textasciicircum{}2 is neither topological nor ledger-reducible, and computes Weyl\textasciicircum{}2=0 on the isotropic layer and O(shear\textasciicircum{}2) on an axisymmetric anisotropic model (no linear term) (5 checks)
 \emph{Bound:} NOT-A-PAPER-CLAIM --- locates L-818's boundary and names the open coupling; F5 reserves the verdict; does NOT close the interacting tower (it names it); coefficients cited (t Hooft-Veltman)
\ \emph{Run:} \texttt{python3 receipts/L819\_shear\_boundary\_coupling/S1\_the\_ledgers\_one\_irreducible\_survivor\_is\_the\_conformal\_coupling\_that\_turns\_on\_at\_the\_shear.py}
\item[\label{rcpt:S1_the_weyl_magnitude_is_the_graviton_c_anomaly_and_the_parity_odd_partner_is_a_theta_term}\texttt{S1\_the\_weyl\_magnitude\_is\_the\_graviton\_c\_anomaly\_and\_the\_parity\_odd\_partner\_is\_a\_theta\_term}]\hfill\textsf{[OK]}\\
\textit{sec:lock} \ (PO-6 OWED \#472: the VALUE of the one shear counterterm (Weyl\textasciicircum{}2, c54.219) is 7/40 --- the graviton's conformal-anomaly c-coefficient (b\_4 in the \{Weyl\textasciicircum{}2,GB,R\textasciicircum{}2\} basis). And its parity-odd partner, the Pontryagin density R\textasciitilde{}R that a CIRCULAR polarisation reveals (nonzero, handedness-odd; a linear mode gives zero), is a topological theta-term at constant coefficient (Jackiw-Pi C-tensor prop nabla(coeff)=0; no D=4 gravitational anomaly), so it does NOT enter the field equations --- the dynamical count stays ONE). 
\emph{Computes:} computes R\textasciitilde{}R on linear vs circular TT waves by a finite-difference Riemann (0 / +4.5 / -4.5, parity-odd), pins the b\_4 Weyl\textasciicircum{}2 coefficient 7/40, and cites the topological/theta-term status (Jackiw-Pi, Alvarez-Gaume-Witten) (4 checks)
 \emph{Bound:} NOT-A-PAPER-CLAIM --- OWED \#472, one value + one status on the fixed-background shear; the parity-odd loop coefficient is zero for the pure graviton in D=4; interacting tower untouched
\ \emph{Run:} \texttt{python3 receipts/L821\_weyl\_coefficient\_value/S1\_the\_weyl\_magnitude\_is\_the\_graviton\_c\_anomaly\_and\_the\_parity\_odd\_partner\_is\_a\_theta\_term.py}
\item[\label{rcpt:S2_c1_and_c2_are_one_and_the_source_answers_more}\texttt{S2\_c1\_and\_c2\_are\_one\_and\_the\_source\_answers\_more}]\hfill\textsf{[OK]}\\
\textit{sec:euclidean} \ (C1 and C2 are ONE condition joined by so --- the list is six not seven --- and the same P10 passage states that the free TT sector is a harmonic oscillator whose Hamiltonian is bounded below, which PO-6 asks for). 
\emph{Computes:} asserts the joined sentence verbatim, confirms PO-6s row asks whether it is bounded below and does not carry the answer, and quotes the confinement clause (11 checks)
 \emph{Bound:} NOT-A-PAPER-CLAIM --- corrects a count and surfaces a partial answer; joint satisfiability untested
\ \emph{Run:} \texttt{python3 receipts/L165\_defining\_the\_sum/S2\_c1\_and\_c2\_are\_one\_and\_the\_source\_answers\_more.py}
\item[\label{rcpt:S3_c6_is_a_theorem_not_a_condition}\texttt{S3\_c6\_is\_a\_theorem\_not\_a\_condition}]\hfill\textsf{[OK]}\\
\textit{sec:lock} \ (the C6/C7 tension cannot arise: the per-fibre structure is DERIVED from a commutation relation, not a requirement on a measure --- so the conditions list is five plus one theorem). 
\emph{Computes:} asserts D1s commutation, direct-integral and fibre-by-fibre statements verbatim, and S1s misclassification of C6 as a condition (8 checks)
 \emph{Bound:} NOT-A-PAPER-CLAIM --- reclassifies a list entry; joint satisfiability untested
\ \emph{Run:} \texttt{python3 receipts/L165\_defining\_the\_sum/S3\_c6\_is\_a\_theorem\_not\_a\_condition.py}
\item[\label{rcpt:S4_the_open_half_is_the_floor}\texttt{S4\_the\_open\_half\_is\_the\_floor}]\hfill\textsf{[OK]}\\
\textit{sec:lock} \ (PO-6s open half is the FLOOR: the direct-integral decomposition needs only that both sides of the threshold are occupied, and P10 leaves whether the complete Gamma-hat is bounded below open in its own voice). 
\emph{Computes:} asserts P10s uses-only clause, its leaves-open clause, the sum-of-squares form and the structural commutation verbatim, and the PO-6 rows own wording (7 checks)
 \emph{Bound:} NOT-A-PAPER-CLAIM --- narrows a register row and corrects r2611s caveat
\ \emph{Run:} \texttt{python3 receipts/L165\_defining\_the\_sum/S4\_the\_open\_half\_is\_the\_floor.py}
\item[\label{rcpt:D3_the_floor_does_not_survive_the_cubic}\texttt{D3\_the\_floor\_does\_not\_survive\_the\_cubic}]\hfill\textsf{[OK]}\\
\textit{sec:lock} \ (PO-6s floor question answered NO at the order the paper names: the cubic enters as pi\textasciicircum{}2 phi, a square times a SIGNED field, so Gamma-hat is unbounded below wherever phi < -c/g3). 
\emph{Computes:} asserts P10s leaves-open clause, its leading-order form and its naming of the cubic verbatim, and computes the sign structure and divergence (8 checks)
 \emph{Bound:} NOT-A-PAPER-CLAIM --- answers a narrowed dark half; the state question is untouched
\ \emph{Run:} \texttt{python3 receipts/L165\_interacting\_tower/D3\_the\_floor\_does\_not\_survive\_the\_cubic.py}
\item[\label{rcpt:D4_the_state_fixer_has_its_own_threshold}\texttt{D4\_the\_state\_fixer\_has\_its\_own\_threshold}]\hfill\textsf{[OK]}\\
\textit{sec:lock} \ (the state-fixing mechanism has its own threshold: nu = sqrt(Gamma + 1/4) is real only for Gamma >= -1/4, below which both exponents go complex and no regular branch exists --- and r2651s operator reaches -47.8). 
\emph{Computes:} asserts P10s thermal-regularity sentence and its common-to-every-fibre clause verbatim, computes nu across Gammas range, and recomputes r2651s minimum (11 checks)
 \emph{Bound:} NOT-A-PAPER-CLAIM --- sharpens PO-6s successor question; whether the region is realised is untouched
\ \emph{Run:} \texttt{python3 receipts/L165\_interacting\_tower/D4\_the\_state\_fixer\_has\_its\_own\_threshold.py}
\item[\label{rcpt:D50_the_floor_does_survive_and_the_paper_said_so_232_revisions_earlier}\texttt{D50\_the\_floor\_does\_survive\_and\_the\_paper\_said\_so\_232\_revisions\_earlier}]\hfill\textsf{[OK]}\\
\textit{sec:deparam} \ ( **`D3` (r2651) HAS THE FLOOR BACKWARDS --- P10 CALLS THE CUBIC'S UNBOUNDEDNESS AN "ARTEFACT OF TRUNCATION" IN THE SAME PARAGRAPH, AND HAS SINCE r2419.** The cubic is the first-order term of \$\textbackslash\{\}pi\textasciicircum{}2/(1+\textbackslash\{\}lambda\textbackslash\{\}phi)\$, whose full coefficient is positive on non-degenerate metrics --- so the floor survives and r2652's sub-\$(-1/4)\$ question has no object. Lead `L-542`; VEIN `L-165`). 
\emph{Computes:} verifies termwise to eighth order that the cubic IS the expansion's first-order term; samples both forms on 200,000 identical points (truncation \$-892.6\$, resummed never below \$\textbackslash\{\}gamma\$); six source checks in P10 including the affirmative verdict and its interior-only SCOPE; dates the clause by reading r2419's own copy of the file; and re-runs `P10\_gamma\_hat\_is\_bounded\_below`, which passes
 \emph{Bound:} LATENT --- no new physics; the answer was in the paper. NOT claimed: that `PO-6` closes (the ultraviolet definition of the tower sums is open in the same sentence), or that the floor is proved for the full interacting theory.  `D3`/`D4` and the `PO-6` row are NOT touched --- 56's band and a protected row; routed
\ \emph{Run:} \texttt{python3 receipts/L165\_interacting\_tower/D50\_the\_floor\_does\_survive\_and\_the\_paper\_said\_so\_232\_revisions\_earlier.py}
\item[\label{rcpt:S50_the_counterterm_basis_is_one_dimensional_because_the_background_family_is}\texttt{S50\_the\_counterterm\_basis\_is\_one\_dimensional\_because\_the\_background\_family\_is}]\hfill\textsf{[OK]}\\
\textit{sec:deparam} \ (**`PO-6`'s UV HALF: THE COUNTERTERM BASIS IS ONE-DIMENSIONAL BECAUSE THE ADMITTED BACKGROUND FAMILY IS.** Every curvature invariant on a maximally symmetric background is a pure power of \$1/\textbackslash\{\}alpha\textasciicircum{}2\$, and the substrates form a one-parameter family, so a divergence of any degree needs exactly one counterterm --- which the framework carries.  Scoped: the tower lives on the layer, whose \$R\$ runs. Lead `L-543`; VEIN `L-165`). 
\emph{Computes:}  **AMENDED c54.211 (`L-544`): c54.210's SCOPE IS WITHDRAWN.** It put P15's observable rate where P10's slicing belongs; P10's own \$a(T)=\textbackslash\{\}alpha\textbackslash\{\}cosh(T/\textbackslash\{\}alpha)\$ has \$R=12/\textbackslash\{\}alpha\textasciicircum{}2\$ CONSTANT, so the degeneracy holds on the tower's own background and the result is STRONGER than claimed. The real limit is the coupled sector, where the scale factor is quantized and there is no fixed background. Both geometries computed side by side so the error stays checkable. --- computes the quadratic invariants at \$D=4,5,6\$ and pins their \$\textbackslash\{\}alpha\$-exponent at exactly 4; five standing-sentence checks (p0's invariant fact, p0's foliation, `L-533`'s scale-only descent); computes the \$\textbackslash\{\}sinh\textasciicircum{}\{2/3\}\$ layer's Ricci scalar and asserts it NON-constant with a \$12H\textasciicircum{}2\$ late limit; and counts `counterterm` in P10 from git HEAD (0) against now (7)
 \emph{Bound:} LATENT, and an ASSEMBLY not a discovery. NOT claimed: that `PO-6` closes, that the sub-leading divergences vanish, that this is renormalisability, or that the layer breaks it --- **that is the question**
\ \emph{Run:} \texttt{python3 receipts/L165\_defining\_the\_sum/S50\_the\_counterterm\_basis\_is\_one\_dimensional\_because\_the\_background\_family\_is.py}
\item[\label{rcpt:S7_the_successor_was_misaimed}\texttt{S7\_the\_successor\_was\_misaimed}]\hfill\textsf{[OK]}\\
\textit{sec:tower} \ (L-543 INHERITS r2677s WITHDRAWN OBJECT: it asks about a RUNNING background and the free tower sits on P10s own slicing at R = 12/alpha\textasciicircum{}2, CONSTANT --- and the layers late limit IS that value, so it asymptotes to it). 
\emph{Computes:} computes both Ricci scalars, shows the layers late limit equals P10s constant under H -> 1/alpha, and pins the successors wording and the back-reaction clause from the row (6 checks)
 \emph{Bound:} NOT-A-PAPER-CLAIM --- withdraws a successor this line registered; the back-reaction question stands
\ \emph{Run:} \texttt{python3 receipts/L165\_interacting\_tower/S7\_the\_successor\_was\_misaimed.py}
\item[\label{rcpt:S8_the_floor_follows}\texttt{S8\_the\_floor\_follows}]\hfill\textsf{[OK]}\\
\textit{sec:lock} \ (PO-6 HAS BEEN CARRYING THE WRONG QUESTION: the complete Gamma-hat inherits its floor from the same non-degeneracy r2671 established, so boundedness is answered YES --- what remains is WHERE the spectrum sits relative to 3/4). 
\emph{Computes:} asserts P10s promotion, inverse-square and Weyl-alternative clauses verbatim, quotes the straddle receipts does-not-need-a-lower-bound line, and shows the truncated coefficient goes negative where the full one stays positive (7 checks)
 \emph{Bound:} NOT-A-PAPER-CLAIM --- the spectrum is untouched; P10 explicitly declines the floor and the decomposition survives without it
\ \emph{Run:} \texttt{python3 receipts/L165\_interacting\_tower/S8\_the\_floor\_follows.py}
\item[\label{rcpt:S9_the_ordering_decides_it}\texttt{S9\_the\_ordering\_decides\_it}]\hfill\textsf{[OK]}\\
\textit{sec:lock} \ (PO-6s SPECTRAL QUESTION IS DECIDED BY OPERATOR ORDERING, which P10 never names: normal-ordered puts the vacuum at 1/4 (BELOW 3/4, freedom survives), symmetric puts N=1 exactly AT 3/4 --- and 3/4 = 1/4 + 1/2, the free coefficient plus one zero-point quantum). 
\emph{Computes:} computes both ordering branches against the threshold, verifies the 1/4+1/2 identity, counts normal-order and symmetric-order at zero in the paper, and pins the both-sides-occupied clause (9 checks)
 \emph{Bound:} NOT-A-PAPER-CLAIM --- neither ordering is asserted correct; P10s decomposition survives either
\ \emph{Run:} \texttt{python3 receipts/L165\_interacting\_tower/S9\_the\_ordering\_decides\_it.py}
\item[\label{rcpt:Q1_the_degeneracy_is_conformal_flatness_not_maximal_symmetry_so_no_scale_factor_can_break_it}\texttt{Q1\_the\_degeneracy\_is\_conformal\_flatness\_not\_maximal\_symmetry\_so\_no\_scale\_factor\_can\_break\_it}]\hfill\textsf{[OK]}\\
\textit{sec:tower} \ (PO-6's DARK HALF HAS AN OBJECT: r2677 stated TWO degeneracies as one. The quadratic trio collapses because the deficit is exactly Weyl\textasciicircum{}2 and EVERY FRW is conformally flat for EVERY a(T) -- not because of maximal symmetry -- so no scale factor can break it and back-reaction cannot; the different-dimension half does fail. The real limit is the SHEAR). 
\emph{Computes:} R, Ric\textasciicircum{}2 and Riem\textasciicircum{}2 computed from the FRW metric with a(T) a free function at k = +1, 0, -1 (Weyl\textasciicircum{}2 identically zero); the Gauss-Bonnet density shown an exact total derivative, sqrt(g)(R\textasciicircum{}2-4Ric\textasciicircum{}2+Riem\textasciicircum{}2) = d/dT[24(a'\textasciicircum{}3/3+a')]; the two relations solved to give int Ric\textasciicircum{}2 = A/3 - chi/2 and int Riem\textasciicircum{}2 = A/3 - chi; S50's D=4 row reproduced from the metric (144/36/24); P15's running layer recomputed from its own metric with the deficit zero by a second route in u = exp(3Ht/2); and Bianchi I giving Weyl\textasciicircum{}2 = 4 sigma\textasciicircum{}2 + O(sigma\textasciicircum{}4). CONTROLS: SdS gives Weyl\textasciicircum{}2 = 48M\textasciicircum{}2/r\textasciicircum{}6 so the identity is not vacuous; the layer's R\textasciicircum{}2 is t-dependent so the different-dimension degeneracy really does fail; and the numeric residual is a PRECISION-SCALING test (3.06e-37 at 40 digits, 2.24e-77 at 80) rather than a fitted threshold
 \emph{Bound:} not that the coupled sector is renormalizable -- this says which functionals a divergence can NEED, not that they are absorbable; no heat-kernel coefficient computed; Bianchi I is a HOMOGENEOUS shear, fixing the ORDER at which conformal flatness fails and not the mode-by-mode statement on the tower; and not a closure -- PO-6 stays open
\ \emph{Run:} \texttt{python3 receipts/L549\_coupled\_counterterms/Q1\_the\_degeneracy\_is\_conformal\_flatness\_not\_maximal\_symmetry\_so\_no\_scale\_factor\_can\_break\_it.py}
\item[\label{rcpt:S10_the_halves_meet_at_the_shear}\texttt{S10\_the\_halves\_meet\_at\_the\_shear}]\hfill\textsf{[OK]}\\
\textit{sec:lock} \ (PO-6s TWO DECLARED HALVES MEET AT THE SHEAR: the counterterm half ends where conformal flatness ends (C\textasciicircum{}2 = 4 sigma\textasciicircum{}2 + O(sigma\textasciicircum{}4)) and P10 names the tower as the TRANSVERSE-TRACELESS graviton modes --- which ARE the shear). 
\emph{Computes:} checks C\textasciicircum{}2 vanishes at sigma=0, is quadratic-leading and non-zero away from it, confirms the source is trace-free, and quotes P10s tower identification twice (6 checks)
 \emph{Bound:} NOT-A-PAPER-CLAIM --- the second-order basis is not computed; cc54 derived the coefficient and its structure is checked here
\ \emph{Run:} \texttt{python3 receipts/L165\_interacting\_tower/S10\_the\_halves\_meet\_at\_the\_shear.py}
\item[\label{rcpt:S11_the_ordering_question_dissolves}\texttt{S11\_the\_ordering\_question\_dissolves}]\hfill\textsf{[OK]}\\
\textit{sec:lock} \ (PO-6s ORDERING QUESTION DISSOLVES: P10 gives deficiency indices (1,1) INDEPENDENTLY of ordering with the coefficient <= 1/4 across the natural family, closes the residual via the horizons Hartle-Hawking state, and answers the coupled straddle fibre by fibre). 
\emph{Computes:} pins all six of P10s own statements --- the independence, the family bound, the non-uniqueness, the Friedrichs selection, the fibre-by-fibre condition, and what remains open (6 checks)
 \emph{Bound:} NOT-A-PAPER-CLAIM --- the interacting tower is still open; the correction is to WHICH problem the row owes
\ \emph{Run:} \texttt{python3 receipts/L165\_interacting\_tower/S11\_the\_ordering\_question\_dissolves.py}
\item[\label{rcpt:P10_ordering_selection_is_external}\texttt{P10\_ordering\_selection\_is\_external}]\hfill\textsf{[OK]}\\
\textit{sec:lock / P7 frontier:quantum} \ (**PO-15 THE ORDERING EXHAUSTION: nothing internal to the construction selects the operator ordering (min Gamma-hat = 1/4 normal vs 3/4 symmetric), and the reason is that selecting it would be solving the cosmological-constant problem.** Completes S11 (boundary CLOSURE is ordering-independent) and the r3014 thermal-state elimination (the thermal state acts downstream). ** THE PIVOT: ** the ordering is BULK-INERT (normal vs symmetric differ by the c-number zero-point 1/2 hbar omega\_n, a global phase under the deparametrized evolution) and BOUNDARY-PHYSICAL (the difference surfaces only in the horizon coefficient, 1/4 -> 3/4 = 1/4 + 1/2, exactly one zero-point quantum), so "which ordering" IS "does the tower's zero-point energy gravitate at the horizon". ** THE EXHAUSTION: ** (1) the thermal state and (3) the seam's characteristic structure act DOWNSTREAM (they close the extension GIVEN Gamma-hat; a regular branch x\textasciicircum{}(1/2+nu) exists for BOTH orderings); (2) the substrate isometry, (5) positivity, (6) the single-scale ledger and (7) covariance under configuration redefinition each respect BOTH orderings; (4) the deparametrization makes normal ordering AVAILABLE (a preferred vacuum) but not mandatory, and by SOLVING the constraint removes the anomaly-freedom lever a Wheeler-DeWitt quantization would have used. The two external naturalness heuristics DISAGREE (continuity-with-free -> 1/4, predictivity -> 3/4). ** RESULT: ** the choice is EXTERNAL because it IS the cc problem reached from inside -- an epistemic gap (one datum: does vacuum energy gravitate), NOT an ontological family, so admissible under the epistemic reading.). 
\emph{Computes:} the ordering difference asserted = the zero-point c-number 1/2 hbar omega; the c-number-is-a-global-phase fact asserted numerically (identical amplitude moduli to 1e-12, ratio = e\textasciicircum{}-ict); the boundary shift 1/4 -> 3/4 asserted to be exactly one zero-point quantum with 3/4 the threshold; the regular branch nu=sqrt(Gamma+1/4) asserted real for BOTH orderings (thermal closure ordering-independent); both floors asserted finite; the seven-candidate enumeration stated with a structural disqualification each and no survivor
 \emph{Bound:} PO-15 answered as the row asked -- the choice is genuinely external, the enumeration is the evidence, the externality is the cc problem localized. NO ordering is picked and none is claimed derivable; the interacting-tower definition (PO-6) stays separately open; builds on S11/r3014, not a re-derivation of the Friedrichs selection
\ \emph{Run:} \texttt{python3 receipts/P10\_canonical\_time/P10\_ordering\_selection\_is\_external.py}
\item[\label{rcpt:S12_the_boundaries_coincide}\texttt{S12\_the\_boundaries\_coincide}]\hfill\textsf{[OK]}\\
\textit{sec:lock} \ (cc54s L-818 BOUNDARY, r2743s SHEAR FINDING AND P10s OPEN ITEM ARE ONE BOUNDARY: an FRW layer is shear-free by construction, the degeneracy ends at the shear, and P10 names the transverse-traceless tower as what remains open). 
\emph{Computes:} verifies cc54s Step 1 (Weyl\textasciicircum{}2 = 0 with R running, all three k) and pins P10s two namings of the tower (5 checks)
 \emph{Bound:} NOT-A-PAPER-CLAIM --- cc54s five sympy steps are not rerun; the coincidence is by identification, each link a corpus statement
\ \emph{Run:} \texttt{python3 receipts/L165\_interacting\_tower/S12\_the\_boundaries\_coincide.py}
\item[\label{rcpt:S13_the_ledger_does_not_reach_second_order}\texttt{S13\_the\_ledger\_does\_not\_reach\_second\_order}]\hfill\textsf{[OK]}\\
\textit{sec:lock} \ (AT SECOND ORDER IN THE SHEAR THE QUADRATIC BASIS IS GENUINELY THREE-DIMENSIONAL: GB and C\textasciicircum{}2 span rank 2 of 3, C\textasciicircum{}2 vanishes identically shear-free so L-818 never had to route it, and it has NONE of L-818s three routes). 
\emph{Computes:} computes the rank and nullspace of \{GB, C\textasciicircum{}2\}, confirms C\textasciicircum{}2 = 0 at all three k on the shear-free layer, and pins p0s one-constant claim as a statement about the FACES (6 checks)
 \emph{Bound:} BOUNDED NEGATIVE --- one order of one expansion; not a verdict on the construction, and the coefficient is not computed
\ \emph{Run:} \texttt{python3 receipts/L165\_interacting\_tower/S13\_the\_ledger\_does\_not\_reach\_second\_order.py}
\item[\label{rcpt:S1_the_shear_needs_exactly_one_new_counterterm_and_my_own_count_was_one_too_many}\texttt{S1\_the\_shear\_needs\_exactly\_one\_new\_counterterm\_and\_my\_own\_count\_was\_one\_too\_many}]\hfill\textsf{[OK]}\\
\textit{sec:tower} \ (PO-6's OWED SHEAR CALCULATION RUN: on a propagating TT mode C\textasciicircum{}2 is O(h\textasciicircum{}2) and carries FOUR derivatives, but delta\textasciicircum{}(1)R = 0 makes int sqrt(g) R\textasciicircum{}2 a shift of Einstein-Hilbert -- so the shear costs exactly ONE new counterterm, Weyl-squared, and c54.215's working count of two was one too many). 
\emph{Computes:} R, Ric\textasciicircum{}2, Riem\textasciicircum{}2 and C\textasciicircum{}2 built from the corpus's own L801/N1 TT-wave-on-de-Sitter metric and expanded to O(eps\textasciicircum{}2); the derivative content fixed by freezing the oscillatory factors and taking the total degree in (H,k,omega), giving 4 against sigma\textasciicircum{}2's 2; delta\textasciicircum{}(1)R shown to vanish and det g shown eps-independent, hence R\textasciicircum{}2 at second order equals 2 Rbar times R at second order, pointwise; the span solved from Gauss-Bonnet plus the Weyl definition, giving B and C both in span of A and C\textasciicircum{}2; and the Pontryagin density computed for linear, circular and opposite-handed polarisation by a finite-difference pipeline cross-checked against a symbolic evaluation to six digits. CONTROLS: the linear mode returns zero for the parity-odd invariant, which is what catches the incomplete basis; the handedness flips its sign; and at eps = 0 everything reduces to c54.215's one dimension
 \emph{Bound:} no heat-kernel COEFFICIENT is computed -- this says which functionals a divergence can need, not with what coefficient; the mode is a plane-wave TT on flat-sliced de Sitter and not P10's S\textasciicircum{}3 harmonics, so the COUNT is mode-independent while the COEFFICIENT is not; the Pontryagin term's total-derivative status is cited, not computed; its consequences for the corpus's chirality result are not raised; and not a closure -- PO-6 stays open
\ \emph{Run:} \texttt{python3 receipts/L553\_the\_shear\_counterterm/S1\_the\_shear\_needs\_exactly\_one\_new\_counterterm\_and\_my\_own\_count\_was\_one\_too\_many.py}
\item[\label{rcpt:D1_the_degeneracy_carrying_the_quartic_was_never_derived_and_its_constant_is_the_component_count}\texttt{D1\_the\_degeneracy\_carrying\_the\_quartic\_was\_never\_derived\_and\_its\_constant\_is\_the\_component\_count}]\hfill\textsf{[OK]}\\
\textit{sec:tower} \ (THE DEGENERACY CARRYING PO-6's QUARTIC WAS NEVER DERIVED: asserted in P10 and asserted-unchecked in D2. Derived from Peter-Weyl on the parallelizable S\textasciicircum{}3 it is g(n) = 2(n-1)(n+3), n >= 2, so the shell contribution is 2n\textasciicircum{}3 -- and the constant is the propagating-component count, which makes D2's "any tensor rank" false). 
\emph{Computes:} the level-j totals of frame-valued fields on S\textasciicircum{}3 = SU(2) for frame-spin 0, 1, 2, coming out as 1, 3, 5 times (2j+1)\textasciicircum{}2 -- the component counts, which is the check the decomposition had to pass; the TT part as the two extreme summands, 2(2j+1)\textasciicircum{}2, verified to be exactly 2/5 of the symmetric-tracefree total; the closed form by eigenvalue with its floor g(2) = 10 matching P10's independently-asserted n >= 2; the cumulative sum and its leading coefficient 2/3; and Weyl's law calibrated on the scalar tower, whose degeneracy follows exactly from Peter-Weyl, with the ratio running 1.01505 to 1.00038. CONTROLS: the scalar tower's leading coefficient 1/3 against the tensor tower's 2/3, which is what refutes "any tensor rank"; and a REPORTED NEGATIVE -- matching Weyl to subleading order fails to discriminate the candidate closed forms, because on a closed curved manifold that term carries the curvature through the heat kernel
 \emph{Bound:} not that the eigenvalue assignment is re-derived, mu\_n\textasciicircum{}2 = n(n+2)-2 being P10's and taken as given; no heat-kernel coefficient, still none computed anywhere; this is the FREE spectrum's mode count and not the interacting tower; and not a closure -- PO-6 stays open, this replaces an unchecked input with a derived one
\ \emph{Run:} \texttt{python3 receipts/L554\_tower\_degeneracy/D1\_the\_degeneracy\_carrying\_the\_quartic\_was\_never\_derived\_and\_its\_constant\_is\_the\_component\_count.py}
\end{description}\endgroup
